\documentclass[letterpaper,11pt]{moderncv}
\moderncvstyle{banking}
\usepackage{xpatch}
\moderncvcolor{black}
\AtBeginDocument{\definecolor{color2}{rgb}{0,0,0}}
\usepackage[letterpaper, margin=0.6in]{geometry}
\usepackage{xcolor}
\usepackage{graphicx}
\usepackage{fontspec}
\setmainfont{Times New Roman}
\pagenumbering{arabic}
\usepackage[unicode]{hyperref}
\sloppy
\usepackage{enumitem}

\usepackage[english]{babel}
\usepackage[autostyle, english = american]{csquotes}
\MakeOuterQuote{"}

%==================== PERSONAL INFO ====================%
\name{\textcolor{black}{Saeyoung}}{Kim}
\address{Cambridge, MA, USA}
\phone[mobile]{+1 (351) 322 5510}
\email{saeyoung\_kim@uml.edu}
\social[linkedin]{saeyoung-macx-kim}
\social[github]{PrimaryAnomaly}

\recipient{Hiring Team}{Fabri\\Boston, MA}
\date{\today}
\opening{Dear Hiring Team,}
\closing{Sincerely,}

\begin{document}

\makelettertitle

I'm applying for the Principal Mechanical Systems Engineer position. My background is a mix of mechanical engineering fundamentals and years of hands-on hardware development, and while my recent work has been closer to the electrical side, the core of what I do is the same thing this role asks for: taking a problem, working through the physics, designing and building hardware, testing it, and iterating until it works.

My foundation is in mechanical engineering. My undergrad was in MechE, where I learned to think in terms of mechanics, materials, and manufacturing processes. That foundation carried through my PhD, where I designed, fabricated, and tested multi-sensor hardware systems from scratch. The work involved CAD modeling, cleanroom microfabrication (thin-film deposition, photolithography, chemical etching), precision mechanical assembly, and building custom test fixtures. Every revision was a full design-build-test-fail-learn cycle, often under tight timelines with limited resources. I managed the \$220K grant that funded the hardware program and made all the build-vs-buy decisions along the way.

At Hanwha Systems, I worked at the system level on military satellite programs, which is where I developed the systems engineering mindset this role requires. I tracked linked requirements across electrical, mechanical, and software interfaces, coordinated across disciplines, and designed test infrastructure that had to work reliably under real constraints. I also contributed to facility architecture planning for a new production line, thinking about throughput, footprint, and scalability.

I'm not going to pretend I have deep turbomachinery or casting equipment experience. What I bring is strong physics fundamentals across multiple domains, a track record of taking hardware from concept to working prototype fast, and the kind of ownership mentality that startup environments require. I learn new domains quickly and I'm not afraid to be on the shop floor.

On ITAR: my green card is currently pending, and I expect to have it in hand within the next few months. I'm happy to discuss timing directly. I'm based in Cambridge, so Boston Metrowest is an easy commute.

\makeletterclosing

\newpage

\makecvtitle
\vspace{-3em}

%==================== SUMMARY ====================%
\section{Professional Summary}
Systems engineer with a mechanical engineering foundation and hands-on experience designing, building, and testing hardware systems from concept through qualification. Background spans CAD modeling, precision fabrication, mechanical assembly, test fixture development, and cross-disciplinary systems integration for aerospace programs. Comfortable working across mechanical, electrical, thermal, and process domains. Experienced managing hardware programs end-to-end, including build-vs-buy decisions, vendor coordination, and rapid prototyping under tight timelines.

\section{Technical Skills}
\begin{minipage}[t]{0.48\textwidth}
\cvitem{Design}{CAD (Fusion360, SolidWorks, AutoCAD), Mechanical Assembly, DFM, Test Fixture Design}
\cvitem{Fabrication}{Cleanroom Processing (Thin-Film Deposition, Photolithography, Etching), 3D Printing, CNC, Soldering, PCB Assembly}
\cvitem{Programming}{Python, C/C++, MATLAB, Bash}
\end{minipage}%
\hfill
\begin{minipage}[t]{0.48\textwidth}
\cvitem{Systems}{Requirements Tracking, Interface Management, V\&V, System-Level Integration}
\cvitem{Instruments}{Oscilloscopes, DMMs, Power Supplies, Function Generators, Calipers, Custom Test Fixtures}
\cvitem{Analysis}{Root Cause Analysis, Process Monitoring, Data-Driven Diagnostics, SPC}
\end{minipage}

%==================== EXPERIENCE ====================%
\section{Professional Experience}

\cventry{2025.07--Present}{\textbf{Postdoctoral Associate} | Process Monitoring and Data-Driven Diagnostics}{University of Massachusetts Lowell}{Lowell, MA, USA}{}{
\begin{itemize}[leftmargin=1cm, itemsep=0pt, parsep=0pt]
    \item Built Python-based diagnostic software for real-time manufacturing process monitoring across 100+ production batches
    \item Developed automated anomaly detection and quality prediction pipelines for multi-day batch manufacturing cycles
    \item Translated cross-functional requirements into standardized procedures and technical documentation
\end{itemize}}

\cventry{2024.02--2025.05}{\textbf{Systems Integration Engineer} | Satellite Assembly, Integration, and Testing}{Hanwha Systems}{Seoul, South Korea}{}{
\begin{itemize}[leftmargin=1cm, itemsep=0pt, parsep=0pt]
    \item Worked as a systems engineer across electrical, mechanical, and software interfaces on military SAR satellite programs
    \item Designed and commissioned test equipment architecture for subsystem and system-level hardware validation
    \item Tracked linked requirements and interface dependencies across tightly coupled disciplines during EM, QM, and FM builds
    \item Performed hands-on mechanical assembly, electrical integration, and troubleshooting in cleanroom environments
    \item Contributed to facility architecture planning for a high-throughput satellite production line, addressing scalability and footprint
    \item Managed purchasing, integration, and commissioning of functional test systems
    \item Coordinated across engineering teams and program leadership to resolve cross-disciplinary hardware issues on tight schedules
\end{itemize}}

%==================== EDUCATION ====================%
\section{Education}
\cventry{2017.03--2024.02}{\textbf{PhD in Electrical Engineering} | Hardware Systems Development}{Korea University}{Seoul, South Korea}{}{
\begin{itemize}[leftmargin=1cm, itemsep=0pt, parsep=0pt]
    \item Secured and managed \$220K grant for CubeSat payload hardware development, owning all build-vs-buy decisions
    \item Designed and built multi-sensor hardware systems through full cycle: CAD modeling, circuit/PCB design, cleanroom microfabrication, mechanical assembly, and system-level qualification
    \item Performed cleanroom fabrication including thin-film deposition, photolithography, wet/dry chemical etching, and precision mechanical assembly
    \item Built custom test fixtures and data acquisition setups for multi-channel sensor characterization
    \item Executed rapid design-build-test iteration cycles across multiple hardware revisions under constrained timelines
    \item Diagnosed cross-domain failures (mechanical, electrical, firmware) and drove fixes through multiple revision cycles
    \item Published 3 first-author and 10+ co-authored papers on sensor hardware, space payloads, and embedded systems
\end{itemize}}

\cventry{2013.03--2017.02}{\textbf{BSc in Mechanical Engineering}}{Konkuk University}{Seoul, South Korea}{}{
\begin{itemize}[leftmargin=1cm, itemsep=0pt, parsep=0pt]
    \item Coursework in mechanics, dynamics, thermodynamics, fluid mechanics, materials science, and machine design
    \item Hands-on experience in CAD modeling, mechanical prototyping, precision measurement, and microfabrication
\end{itemize}}

%==================== PUBLICATIONS ====================%
\section{Publications}
\subsection{Journal Publications}
\begin{sloppypar}
\begin{itemize}[leftmargin=1cm, itemsep=0pt, parsep=0pt]
    \item \textbf{\underline{Saeyoung Kim}}, Jaewon Lee, Kimia Babaei, Seonkgyu Yoon, "Digital Twins in Lyophilization: Advances in Pharmaceutical Process Development and Optimization," \textit{Biotechnology Advances}, 2026. (Manuscript in Preparation, Target Submission: Q1 2026)
    \item Jinhwee Kil, \textbf{\underline{Saeyoung Kim}}, James Jungho Pak, "Investigation of thermal management strategies for biological CubeSat payloads," Acta Astronautica, 2024, Art. no. 228. DOI: 10.1016/j.actaastro.2024.11.059.
    \item \textbf{\underline{Kim, S.}}; Park, S.; Pak, J.J., "Multi-Modal Multi-Array Electrochemical and Optical Sensor Suite for a Biological CubeSat Payload," MDPI Sensors, 2024, 24, Art. no. 265. DOI: 10.3390/s24010265.
    \item Ji-Hoon Han, Sang Hyun Park, \textbf{\underline{Saeyoung Kim}}, James Jungho Pak, "A Performance Improvement of Enzyme-Based Electrochemical Lactate Sensor Fabricated by Electroplating Novel PdCu Mediator on a Laser Induced Graphene Electrode," Bioelectrochemistry, 2022, 148, 108259. DOI: 10.1016/j.bioelechem.2022.108259.
    \item \textbf{\underline{Saeyoung Kim}}, Ji-Hoon Han, Beelee Chua, James Jungho Pak, "A pH Sensing Pipette for Cross-Contamination Prevention in Industrial Fermentation," IEEE Transactions on Industrial Electronics, 2022, 69(7), 7461-7469. DOI: 10.1109/TIE.2021.3099251.
    \item S.D. Raut, N.M. Shinde, B.G. Ghule, \textbf{\underline{S. Kim}}, J.J. Pak, Q. Xia, R.S. Mane, "Room-Temperature Solution Processed Sharp-Edged Nanoshapes of," Chemical Engineering Journal, 2022, 433(2), 133627. DOI: 10.1016/j.cej.2021.133627.
    \item S.E. Jeong, \textbf{\underline{S. Kim}}, J.H. Han, J. Jungho Pak, "Simple Laser-Induced Graphene Fiber Electrode Fabrication for High-Performance Heavy-Metal Sensing," Microchemical Journal, 2022, 172, 106950. DOI: 10.1016/j.microc.2021.106950.
    \item Sanghoon Cho, Jungmo Jung, \textbf{\underline{Saeyoung Kim}}, James Pak, "Conduction Mechanism and Synaptic Behaviour of Interfacial Switching AlOσ-based RRAM," Semiconductor Science and Technology, 2020, 35(8), 085006. DOI: 10.1088/1361-6641/ab8d0e.
    \item Ji-Hoon Han, \textbf{\underline{Saeyoung Kim}}, Jaesung Choi, Sora Kang, Youngmi Kim Pak, James Jungho Pak, "Development of Multi-Well-Based Electrochemical Dissolved Oxygen Sensor Array," Sensors and Actuators B: Chemical, 2020, 306, 127465. DOI: 10.1016/j.snb.2019.127465.
\end{itemize}
\end{sloppypar}

\subsection{Conference Presentations}
\begin{sloppypar}
\begin{itemize}[leftmargin=1cm, itemsep=0pt, parsep=0pt]
    \item \textit{Thermal Management System for a Fluidic Card in a Small Satellite Biological Payload}, Korean Institute of Electrical Engineers (KIEE) Autumn Conference, 2022.
    \item \textit{A Pipette with an Integrated pH Sensor for Liquid Sampling Applications}, Korean Institute of Electrical Engineers (KIEE) Summer Conference, 2019.
    \item \textit{IoT System for Monitoring pH, Dissolved Oxygen, Temperature, and Electrical Conductivity in Liquids}, Korean Institute of Electrical Engineers (KIEE) Summer Conference, 2018.
    \item \textit{Development Trends of Non-Invasive Glucose Sensor Technology}, Korean Institute of Electrical Engineers (KIEE) Summer Conference, 2017.
\end{itemize}
\end{sloppypar}

\subsection{Seminar Presentations}
\begin{itemize}[leftmargin=0.5cm, itemsep=2pt, parsep=1pt]
    \item \textit{Bioprocessing in Space: Challenges and Opportunities for Pharmaceutical Manufacturing Beyond Earth}, Lowell Center for Space Science and Technology (LoCSST), UMass Lowell, 2025.
\end{itemize}

\vfill
\centering{References available upon request.}

\end{document}
